% =============================================================================
% IV–V. Experimental Setup / Results and Discussion
% Status: P2 draft — numbers frozen in paper/experiment_freeze.json
% IMPORTANT: Do NOT mix Table I (16 UAV, 5-seed) with Ablation (4 UAV).
% =============================================================================

\section{Experimental Setup}
\label{sec:setup}

\subsection{Environment}
All experiments use the VMAS continuous multi-agent \texttt{navigation} scenario
(obstacle generalizations use \texttt{navigation\_obstacle}).
Unless noted, training uses $102{,}400$ environment frames, seed $42$ for single
runs, and five seeds $\{42,43,44,45,46\}$ for Table~\ref{tab:main}.
Episode horizon is $H{=}128$ steps; success requires goal proximity below a
fixed threshold ($0.3$ world units).
Communication radius $R_c$ controls graph construction for GAT/DSGF; MAPPO does
not use explicit inter-agent messaging (reported communication cost $=0$).

\subsection{Baselines}
We compare:
\begin{itemize}
  \item \textbf{MAPPO}: CTDE PPO without guidance.
  \item \textbf{GAT-MAPPO}: binary radius graph + GAT guide with alignment reward.
  \item \textbf{Transformer / Full Attention}: dense $O(N^{2})$ attention guide.
  \item \textbf{DSGF v2 (Ours)}: quality-weighted sparse graph + temporal GRU +
        residual policy (\emph{frozen}; no further hyperparameter search).
\end{itemize}
Hyperparameters for 16-UAV comparisons follow
\texttt{configs/baseline16/};
DSGF settings match \texttt{configs/dsgf/dsgf\_v2\_frozen.yaml}.

\subsection{Metrics}
\begin{itemize}
  \item \textbf{Success rate} $S$: fraction of agents within the goal threshold
        at evaluation (200 episodes unless stated).
  \item \textbf{Collision rate}: fraction of episodes with collisions.
  \item \textbf{Communication cost}
        $C=\frac{1}{T}\sum_t |E_t|$ (directed edges per timestep).
  \item \textbf{Efficiency} $\eta=S/C$ (communication Pareto analysis).
  \item \textbf{AULC}: area under the training reward curve (stability proxy).
\end{itemize}

\subsection{Evaluation Protocol}
\begin{enumerate}
  \item \textbf{Main comparison (Table~\ref{tab:main})}: $N{=}16$, 5 seeds.
  \item \textbf{Ablation (Table~\ref{tab:ablation})}: $N{=}4$, seed $42$;
        remove Dynamic Graph / Temporal / Residual independently.
  \item \textbf{Generalization (Table~\ref{tab:gen})}: train at obstacle count
        $4$; zero-shot evaluate $\{0,2,4,6,8\}$.
  \item \textbf{Communication (Table~\ref{tab:comm})}: $N{=}16$,
        $R_c\in\{1,2,\mathrm{full}\}$; retrain GAT/DSGF per radius.
  \item \textbf{Scalability}: DSGF at $N\in\{4,8,16,32\}$.
\end{enumerate}


\section{Results and Discussion}
\label{sec:results}

\subsection{Main Comparison}
\label{sec:main-results}
Table~\ref{tab:main} reports mean$\pm$std success over five seeds at $N{=}16$.
DSGF achieves $4.01\%\pm1.68\%$ versus MAPPO $2.40\%\pm0.90\%$ and GAT
$2.04\%\pm0.90\%$; Transformer remains weakest ($0.70\%\pm0.16\%$).
Paired Wilcoxon / $t$-tests against MAPPO yield $p{=}0.063$ / $0.107$
(borderline; limited seed budget), so we claim a consistent numerical
advantage rather than strong statistical superiority.
Dense attention does not help under this swarm navigation regime, supporting
the sparsity premise of DSGF.

% Auto-include CSV commentary:
% Method Success_mean±std AULC_mean±std
% MAPPO 2.40±0.90  4.9995±0.2206
% GAT   2.04±0.90  4.4421±0.2935
% Trans 0.70±0.16  4.7599±0.2866
% DSGF  4.01±1.68  4.0307±0.9969

\paragraph{Learning stability.}
AULC (normalized reward AUC) remains comparable across methods, while eval
success diverges---consistent with the motivation that \emph{training reward
alone} is a poor proxy for coordination success under long horizons.
GAT exhibits classic over-specialization in smaller-$N$ pilots
(early peak then collapse); residual decoupling in DSGF mitigates this failure
mode (Sec.~\ref{sec:ablation-results}).

\subsection{Ablation Study}
\label{sec:ablation-results}
On the controlled 4-UAV setting (Table~\ref{tab:ablation}), Full DSGF reaches
$9.27\%$ success.
Removing Residual collapses performance to $0.22\%$ ($\sim$42$\times$ drop),
dominating other ablations (w/o Dynamic Graph $0.28\%$; w/o Temporal $0.86\%$).
This isolates residual policy decoupling as the primary mechanism behind
long-horizon robustness, rather than graph or memory alone.

\subsection{Generalization to Obstacle Density}
\label{sec:gen-results}
Policies trained at obstacle count $4$ are zero-shot evaluated at
$\{0,2,4,6,8\}$ (Table~\ref{tab:gen}).
All methods remain in a low absolute success regime, reflecting the difficulty
of 4-UAV navigation with randomized obstacles.
Importantly, we do \emph{not} observe catastrophic degradation for DSGF as
obstacle count increases beyond the training density; MAPPO/GAT remain
competitive on some bins.
We therefore claim \textbf{robust transfer under density shift}, not absolute
superiority on every obstacle level.

\subsection{Communication Efficiency}
\label{sec:comm-results}
Table~\ref{tab:comm} and Fig.~\ref{fig:comm} compare success versus communication
cost under radius constraints ($N{=}16$).
Under \textbf{restricted} radii, DSGF outperforms GAT:
$R_c{=}1$ yields $2.22\%$ vs.\ $1.88\%$;
$R_c{=}2$ yields $4.07\%$ vs.\ $1.28\%$.
Under full connectivity, GAT can achieve higher success ($3.19\%$ vs.\ DSGF
$2.41\%$), so we do \textbf{not} claim unconditional communication dominance.
The intended narrative is:
\begin{quote}
DSGF maintains competitive success under restricted communication budgets,
demonstrating robustness to sparse connectivity.
\end{quote}

\subsection{Scalability}
\label{sec:scale-results}
Fig.~\ref{fig:scale} shows DSGF success for $N\in\{4,8,16,32\}$.
Absolute rates remain low as team size grows (all methods $<5\%$ at $N{=}16$
in Table~\ref{tab:main}), highlighting that large-scale continuous coordination
is still an open challenge.
DSGF remains trainable up to $32$ UAVs with the frozen recipe after spawn-density
fixes, providing a scalable experimental protocol rather than saturating
performance claims.

\subsection{Complexity}
\label{sec:complexity-disc}
As summarized in Table~\ref{tab:complexity}, dense GAT/Transformer attention
scales as $O(N^{2}d)$, whereas DSGF message passing is $O(Nkd)$ with an
additional $O(Nd)$ temporal update.
When $k\ll N$, the guidance stack approaches linear scaling in $N$, aligning
with the empirical communication study under small $R_c$.


\section*{Limitations and Discussion}
\label{sec:limitations}

\begin{enumerate}
  \item \textbf{Absolute success remains low} at $N{=}16$ (single-digit \%).
        The contribution is comparative evidence under a hard continuous MARL
        benchmark, not claimed real-world flight readiness.
  \item \textbf{Seed variance} on Table~\ref{tab:main} is non-negligible
        (DSGF std $1.68\%$); five seeds are modest for strict significance.
  \item \textbf{Full-graph regimes}: GAT may outperform DSGF when communication
        is unconstrained; DSGF's advantage concentrates on sparse budgets.
  \item \textbf{Cross-table mixing}: Ablation (4 UAV) and main results (16 UAV)
        use different scales and must not be compared numerically in text.
\end{enumerate}

% Suggested table/figure labels for the master .tex:
% \label{tab:main}        <- table1_main.csv
% \label{tab:ablation}    <- table2_ablation.csv
% \label{tab:gen}         <- table3_generalization.csv
% \label{tab:comm}        <- table4_communication.csv
% \label{tab:complexity}  <- table5_complexity.tex
% \label{fig:comm}        <- fig6_communication.png
% \label{fig:scale}       <- fig3_scalability.png
% \label{fig:ablation}    <- fig4_ablation.png
% \label{fig:gen}         <- fig5_generalization.png
% \label{fig:curve}       <- fig2_learning_curve_5seed.png
